% ============================================================================
% Fichier  : partie4/P04_C15_niveaux_service_continuite.tex
% Titre    : Niveaux de Service, Disponibilité, Continuité et Fournisseurs
% Auteur   : MINKA MI NGUIDJOÏ Thierry Emmanuel
% Version  : 1.0 -- octobre 2026
% ============================================================================
% Sources : ITIL 4 (sources publiques, notions reformulées) ; COBIT 2019,
% fiches APO09, APO10 et DSS04 (dossier Normes) ; loi n° 2010/012, art. 26
% (texte anglais officiel). Laboratoire associé : GCI-LAB-08 (disponibilité
% mensuelle, pénalités en FCFA, RTO et RPO chez KRIBILOG).
% ============================================================================

\chapter{Niveaux de Service, Disponibilité, Continuité et Fournisseurs}
\label{chap:itil-niveaux}
\chaptermark{Niveaux de service et fournisseurs}

\epigraph{%
    « Un engagement que l'on ne mesure pas n'engage personne. »%
}{Adage des acheteurs de services}

\niveauChapitre{Intermédiaire}{%
    Maîtrise du Chapitre~\ref{chap:itil-incidents} (incidents, priorités et
    délais cibles) et du Chapitre~\ref{chap:itil-svs} (utilité et
    garantie).%
}

\begin{boxobjectifs}
\begin{itemize}
    \item Rédiger un SLA, un OLA et un contrat de sous-traitance cohérents.
    \item Calculer une disponibilité et les pénalités contractuelles en
          FCFA.
    \item Conduire une analyse d'impact sur l'activité et fixer RTO et RPO.
    \item Piloter un fournisseur de services informatiques.
\end{itemize}
\end{boxobjectifs}

\bigskip

% ============================================================================
\section*{Question directrice : que vaut une promesse de 99,5~\%~?}
% ============================================================================

Le contrat entre KRIBILOG SA et son infogérant NETPORT SERVICES tient en
une ligne d'engagement~: « disponibilité des serveurs~: 99,5~\% ».
L'exercice~12.3 a montré que cette promesse ne protège pas KRIBILOG.
Mais combien vaut-elle exactement~? Combien d'heures de panne
autorise-t-elle chaque mois~? Que rembourse-t-elle quand un navire attend
à quai, à 12~millions FCFA de surestaries par jour~? Et que se passe-t-il si
le centre de données de NETPORT SERVICES brûle~?

Ce chapitre répond à ces questions. Il présente la gestion des niveaux de
service (section~\ref{sec:c15-sla}), le calcul de la disponibilité
(section~\ref{sec:c15-dispo}) et des pénalités (section~\ref{sec:c15-penalites}),
la continuité de service (section~\ref{sec:c15-continuite}) et la gestion
des fournisseurs (section~\ref{sec:c15-fournisseurs}). Le laboratoire
GCI-LAB-08 en fait réaliser les calculs sur les données de KRIBILOG.

\separateur

% ============================================================================
\section{Les Engagements de Service}
\label{sec:c15-sla}
% ============================================================================

\subsection{Trois niveaux d'accord}

\begin{boxdef}
\termecle{Accord de niveau de service}\index{accord de niveau de service}
(\textit{service level agreement}, \sigle{SLA})~: accord documenté entre un
fournisseur de services et son client, qui précise le service fourni et
les niveaux de service attendus.

\medskip
\termecle{Accord de niveau opérationnel}\index{accord de niveau
opérationnel} (\textit{operational level agreement}, \sigle{OLA})~: accord
entre le fournisseur et une autre équipe interne de la même organisation,
qui soutient la fourniture du service.

\medskip
\termecle{Contrat de sous-traitance}\index{contrat de sous-traitance}
(\textit{underpinning contract})~: contrat entre le fournisseur et un
tiers extérieur dont il dépend pour fournir le service.
\end{boxdef}
{\small\itshape Notions reformulées d'après ITIL~4~\cite{axelos2019itil4}.}

\noindent\begin{minipage}{\linewidth}
\begin{center}
\begin{tikzpicture}[
    acteur/.style = {
        draw=CouleurBordDef, fill=CouleurDef, rounded corners=4pt,
        text width=2.8cm, minimum height=1.0cm, align=center,
        font=\footnotesize\bfseries\color{CouleurPrimaire}
    },
    accord/.style = {font=\scriptsize\bfseries\color{CouleurAccent}},
    lien/.style = {{Stealth[length=5pt]}-{Stealth[length=5pt]}, thick,
        color=CouleurSecondaire}
]
\node[acteur] (cl) at (0,0)    {Armateur\\ (client)};
\node[acteur] (kr) at (4.6,0)  {KRIBILOG\\ (fournisseur)};
\node[acteur] (in) at (4.6,-2.2) {Équipe réseau\\ interne};
\node[acteur] (np) at (10.0,0) {NETPORT SERVICES\\ (infogérant)};
\draw[lien] (cl) -- node[above, accord]{SLA} (kr);
\draw[lien] (kr) -- node[right, accord]{OLA} (in);
\draw[lien] (kr) -- node[below, accord, align=center]{contrat de\\
    sous-traitance} (np);
\end{tikzpicture}
\end{center}
\vspace{-0.4em}
\begin{center}{\small\itshape Figure~15.1. SLA, OLA et contrat de
sous-traitance du service de suivi des conteneurs}\end{center}
\end{minipage}

\medskip

La règle de cohérence est simple~: \textbf{on ne peut pas promettre au
client plus que ce que garantissent les accords qui soutiennent le
service}. Si KRIBILOG promet à l'armateur un rétablissement en 4~heures,
mais que le contrat de NETPORT SERVICES ne fixe aucun délai
d'intervention, la promesse ne repose sur rien.

\subsection{Rédiger un SLA utile}

\begin{tableaucours}{L{3.2cm}L{4.6cm}L{4.4cm}}{%
    \tableauHeader{Rubrique} &
    \tableauHeader{Contenu} &
    \tableauHeader{Exemple KRIBILOG}
}
Description du service & Ce que le service permet au client
d'obtenir (le résultat) & Localiser en temps réel les conteneurs de
l'armateur \\
Horaires de service & Périodes pendant lesquelles les engagements
s'appliquent & 24~h/24, 7~j/7 \\
Disponibilité & Part du temps de service où le service est utilisable &
99,5~\% par mois \\
Délais de rétablissement & Par priorité (Chapitre~\ref{chap:itil-incidents})
& P1 en 4~heures, P2 en 8~heures \\
Indicateur de résultat & Mesure ce que le client vit & Part des conteneurs
localisés en moins de 5~minutes~: 99~\% \\
Mesure et compte rendu & Source des données, fréquence & Rapport mensuel
produit par KRIBILOG à partir de son outil \\
Pénalités & Conséquences du non-respect & Section~\ref{sec:c15-penalites} \\
Revue & Fréquence de révision de l'accord & Comité de service
trimestriel \\
\end{tableaucours}
\vspace{-0.8em}
\begin{center}{\small\itshape Tableau~15.1. Les rubriques d'un accord de
niveau de service}\end{center}

\begin{boxalert}
\textbf{Violation typique.} Tous les indicateurs du rapport mensuel sont
au vert, mais le client est mécontent. ITIL~4 appelle ce phénomène
l'« effet pastèque »~: vert à l'extérieur, rouge à l'intérieur. Il naît
d'engagements qui mesurent des livrables (serveurs disponibles) au lieu
de résultats (conteneurs localisés).

\smallskip
\textbf{Règle~:} chaque SLA comporte au moins un indicateur de résultat,
défini avec le client, et la satisfaction du client est mesurée à côté
des indicateurs techniques.
\end{boxalert}

\begin{boxobjet}{APO09 · Gérer les accords de service}
\textbf{Finalité~:} s'assurer que les services et les niveaux de service
répondent aux besoins actuels et futurs de l'organisation.

\textbf{Objectifs liés~:} EG01, EG08~; AG05.

\textbf{Pratiques~:} APO09.01 identifier les services~; APO09.02
cataloguer les services~; APO09.03 définir et préparer les accords de
service~; APO09.04 surveiller les niveaux de service et en rendre
compte~; APO09.05 revoir les accords de service et les contrats.

\textbf{RACI (extrait, KRIBILOG)~:} APO09.03~: \emph{A} au directeur des
opérations, \emph{R} au responsable des services informatiques, \emph{C}
à l'armateur et à NETPORT SERVICES.

\textbf{Indicateurs~:} part des services couverts par un SLA, part des
engagements tenus, satisfaction des clients.
\end{boxobjet}
{\small\itshape Pratiques et objectifs liés traduits par nos soins à partir
de COBIT~2019~\cite{isaca2018cobitobjectifs}~; finalité reformulée.}

% ============================================================================
\section{Calculer une Disponibilité}
\label{sec:c15-dispo}
% ============================================================================

\subsection{La formule}

\begin{boxdef}
\termecle{Disponibilité}\index{disponibilité}~: aptitude d'un service à
remplir sa fonction quand on en a besoin. Elle se mesure par
\[
    A = \frac{T_{\text{convenu}} - T_{\text{indisponible}}}{T_{\text{convenu}}}
    \times 100,
\]
où $T_{\text{convenu}}$ est le temps de service convenu sur la période et
$T_{\text{indisponible}}$ la durée cumulée des interruptions pendant ce
temps.
\end{boxdef}

Sur un mois de 30~jours en service 24~h/24, le temps convenu est de
720~heures. Une disponibilité de 99,5~\% autorise donc $0{,}005 \times 720
= 3{,}6$~heures d'interruption par mois~; 99,9~\% en autorisent
43~minutes~; 99,99~\%, un peu plus de 4~minutes.

\begin{tableaucours}{L{3.0cm}L{4.6cm}L{4.6cm}}{%
    \tableauHeader{Disponibilité} &
    \tableauHeader{Interruption par mois (720 h)} &
    \tableauHeader{Interruption par an (8~760 h)}
}
99~\% & 7,2 heures & 87,6 heures \\
99,5~\% & 3,6 heures & 43,8 heures \\
99,9~\% & 43 minutes & 8,8 heures \\
99,99~\% & 4,3 minutes & 53 minutes \\
\end{tableaucours}
\vspace{-0.8em}
\begin{center}{\small\itshape Tableau~15.2. Ce que signifie un taux de
disponibilité}\end{center}

\subsection{Composants en série et en parallèle}

Un service dépend de plusieurs composants. S'ils sont \textbf{en série}
(il les faut tous), la disponibilité du service est le produit de leurs
disponibilités~: $A = A_1 \times A_2 \times \cdots \times A_n$. S'ils sont
\textbf{en parallèle} (un seul suffit, les autres sont des secours), la
disponibilité est $A = 1 - (1-A_1)(1-A_2)\cdots(1-A_n)$.

\begin{boxscenario}{Opération KRIBI. La disponibilité réelle du suivi}
\textbf{Composants en série.} Serveur d'application (99,9~\%), base de
données (99,8~\%), passerelle des balises (99,5~\%), réseau de
l'opérateur mobile (99,0~\%).

\medskip
\textbf{Calcul.} $0{,}999 \times 0{,}998 \times 0{,}995 \times 0{,}990 =
0{,}9821$, soit \textbf{98,21~\%}~: environ 12,9~heures d'interruption
attendues par mois, alors que l'engagement de NETPORT SERVICES sur les
seuls serveurs (99,5~\%) est tenu.

\medskip
\textbf{Avec un second opérateur mobile en secours.} Le réseau passe à
$1 - (0{,}01 \times 0{,}01) = 0{,}9999$, et le service à $0{,}999 \times
0{,}998 \times 0{,}995 \times 0{,}9999 = 0{,}9919$, soit 99,19~\%.

\medskip
\textbf{Analyse} (confiance élevée). Le maillon faible n'est pas le
serveur, sur lequel porte le contrat, mais le réseau mobile, sur lequel
il ne porte pas. Un SLA sur un composant ne garantit pas le service.
\end{boxscenario}

% ============================================================================
\section{Les Pénalités}
\label{sec:c15-penalites}
% ============================================================================

\subsection{Un mécanisme contractuel}

Les pénalités (souvent appelées « crédits de service ») sanctionnent le
non-respect des engagements. Elles ont trois fonctions~: inciter le
fournisseur, compenser partiellement le client, et documenter les
manquements. Elles ne compensent presque jamais le préjudice réel.

\begin{boxscenario}{Opération KRIBI. Le calcul du mois de juin}
\textbf{Clause du nouveau contrat.} Redevance mensuelle~: 25~millions
FCFA. Engagement~: 99,5~\% de disponibilité du service de suivi, mesurée
de bout en bout. Pénalité~: 4~\% de la redevance par tranche complète de
0,1~point en dessous de l'engagement, plafonnée à 30~\%.

\medskip
\textbf{Faits du mois.} Trois interruptions de 2~heures, 4,5~heures et
1,5~heure, soit 8~heures sur 720.

\medskip
\textbf{Calcul.} Disponibilité~: $(720-8)/720 = 98{,}89~\%$. Écart~: $99{,}5
- 98{,}89 = 0{,}61$~point, soit 6~tranches complètes de 0,1~point.
Pénalité~: $6 \times 4~\% = 24~\%$, sous le plafond, soit $0{,}24 \times 25 =
6$~millions FCFA.

\medskip
\textbf{Mise en perspective.} L'interruption de 4,5~heures a immobilisé
un navire une demi-journée~: environ 6~millions FCFA de surestaries à elle
seule. La pénalité du mois couvre une seule interruption sur trois.
\end{boxscenario}

\begin{boiteRemarque}{Faut-il des pénalités plus lourdes~?}
Pas nécessairement. Des pénalités très lourdes renchérissent le contrat,
poussent le fournisseur à contester les mesures et ne remplacent pas un
service fiable. La protection réelle du client vient d'ailleurs~: des
engagements de résultat, une architecture redondante, un plan de
continuité, un droit de résiliation en cas de manquements répétés et une
supervision effective du fournisseur. La pénalité est un signal, pas une
assurance.
\end{boiteRemarque}

% ============================================================================
\section{La Continuité de Service}
\label{sec:c15-continuite}
% ============================================================================

\subsection{L'analyse d'impact sur l'activité}

La continuité vise à maintenir les services à un niveau acceptable en cas
de sinistre majeur~: incendie du centre de données, coupure prolongée
d'électricité, cyberattaque destructrice. Elle commence par une analyse
d'impact sur l'activité.

\begin{boxdef}
\termecle{Analyse d'impact sur l'activité}\index{analyse d'impact sur
l'activité} (\textit{business impact analysis}, \sigle{BIA})~: analyse qui
identifie les activités essentielles, mesure l'effet de leur interruption
dans le temps, et en déduit les objectifs de reprise.

\medskip
\termecle{Durée maximale d'interruption admissible}\index{durée maximale
d'interruption admissible} (\sigle{DMIA})~: durée au-delà de laquelle
l'interruption d'une activité devient inacceptable pour l'organisation.

\medskip
\termecle{RTO}\index{RTO} (\textit{recovery time objective}, durée
maximale de reprise)~: délai cible dans lequel le service doit être
rétabli après un sinistre. Il doit être inférieur à la DMIA.

\medskip
\termecle{RPO}\index{RPO} (\textit{recovery point objective}, perte de
données maximale admissible)~: quantité maximale de données, exprimée en
durée, que l'on accepte de perdre. Un RPO de 15~minutes impose une
réplication ou une sauvegarde au moins toutes les 15~minutes.
\end{boxdef}

\begin{tableaucours}{L{3.6cm}L{2.0cm}L{1.6cm}L{1.6cm}L{3.0cm}}{%
    \tableauHeader{Service} &
    \tableauHeader{DMIA} &
    \tableauHeader{RTO} &
    \tableauHeader{RPO} &
    \tableauHeader{Dispositif}
}
Système central (SANAGA) & 8 heures & 4 h & 15 min & Site de secours
répliqué \\
Paie des fonctionnaires (ANPP) & 48 heures & 24 h & 1 h & Sauvegarde
horaire, site de secours \\
Suivi des conteneurs (KRIBILOG) & 4 heures & 2 h & 15 min & Réplication,
second opérateur \\
Messagerie (KRIBILOG) & 3 jours & 24 h & 24 h & Sauvegarde quotidienne \\
\end{tableaucours}
\vspace{-0.8em}
\begin{center}{\small\itshape Tableau~15.3. Objectifs de reprise de
quelques services du cours}\end{center}

\subsection{Le plan et ses tests}

Le plan de continuité ne vaut que s'il est \emph{testé}. COBIT y consacre
l'objectif DSS04~; la norme ISO~22301 en formule les exigences
(Chapitre~\ref{chap:normes}).

\begin{boxobjet}{DSS04 · Gérer la continuité}
\textbf{Finalité~:} s'adapter rapidement, poursuivre les activités et
maintenir la disponibilité des ressources et de l'information à un niveau
acceptable en cas de perturbation majeure.

\textbf{Objectifs liés~:} EG01, EG02, EG06, EG08~; AG05, AG07.

\textbf{Pratiques~:} DSS04.01 définir la politique, les objectifs et le
périmètre de la continuité~; DSS04.02 maintenir la résilience~; DSS04.03
élaborer et mettre en œuvre la réponse~; DSS04.04 exercer, tester et revoir
les plans~; DSS04.05 revoir, maintenir et améliorer les plans~; DSS04.06
former~; DSS04.07 gérer les sauvegardes~; DSS04.08 conduire une revue après
la reprise.

\textbf{Indicateurs~:} part des services critiques dont le plan a été testé
dans l'année~; écart entre RTO visé et durée de reprise constatée lors des
tests.
\end{boxobjet}
{\small\itshape Finalité, objectifs liés et pratiques résumés et traduits
par nos soins à partir de COBIT~2019~\cite{isaca2018cobitobjectifs}.}

\begin{tableaucours}{L{3.4cm}L{4.6cm}L{4.2cm}}{%
    \tableauHeader{Type de test} &
    \tableauHeader{Déroulement} &
    \tableauHeader{Ce qu'il prouve}
}
Exercice sur table & Les responsables déroulent le plan sur un scénario,
autour d'une table & Que chacun connaît son rôle \\
Simulation & Les équipes exécutent les procédures sans basculer
réellement & Que les procédures sont praticables \\
Bascule réelle & Le service est réellement repris sur le site de secours
& Que le RTO et le RPO sont tenus \\
\end{tableaucours}
\vspace{-0.8em}
\begin{center}{\small\itshape Tableau~15.4. Types de tests de continuité}\end{center}

\begin{boxalert}
\textbf{Violation typique.} La BANQUE SANAGA disposait d'un plan de
continuité rédigé en 2022. Il n'avait jamais été testé, et le site de
secours qu'il désignait hébergeait encore l'ancienne version du système
central. Le jour de la panne, personne n'a pensé à l'ouvrir.

\smallskip
\textbf{Règle~:} un plan de continuité se teste au moins une fois par an
pour chaque service critique, et se met à jour à chaque changement majeur
(Chapitre~\ref{chap:itil-changements}).
\end{boxalert}

% ============================================================================
\section{Piloter les Fournisseurs}
\label{sec:c15-fournisseurs}
% ============================================================================

\subsection{L'externalisation ne transfère pas la responsabilité}

Le Chapitre~\ref{chap:defaillances} l'a montré avec TSB, le
Chapitre~\ref{chap:normes} avec l'ISO/IEC~20000-1~: confier un service à
un tiers ne décharge pas l'organisation de sa responsabilité. Le droit
camerounais le dit aussi.

\begin{boxjuris}
\textbf{\loicam{n°~2010/012}, article~26, alinéa~1} (extrait, version
anglaise officielle)~:

\medskip
\textit{« Operators of information systems shall take every technical and
administrative measure to ensure the security of services offered. To
this end, they shall have standardized systems enabling them to at all
times identify, assess, process or manage any risk relating to the
security of the information systems of the services provided directly or
indirectly. »}

\medskip
\textbf{Portée.} L'obligation de gérer les risques de sécurité porte sur
les services fournis « directement ou indirectement ». Une organisation qui
exploite ses systèmes par l'intermédiaire d'un infogérant reste tenue
d'identifier, d'évaluer et de traiter les risques de sécurité de ces
services~: elle doit donc s'en donner les moyens dans le contrat (droit
d'audit, information sur les incidents, exigences de sécurité).
\end{boxjuris}

\subsection{Le cycle de la relation fournisseur}

\begin{boxobjet}{APO10 · Gérer les fournisseurs}
\textbf{Finalité~:} obtenir des fournisseurs les services convenus, en
maîtrisant leurs risques, et en entretenant des relations efficaces.

\textbf{Objectifs liés~:} EG01, EG08~; AG05.

\textbf{Pratiques~:} APO10.01 identifier et évaluer les relations et les
contrats~; APO10.02 sélectionner les fournisseurs~; APO10.03 gérer les
relations et les contrats~; APO10.04 gérer les risques liés aux
fournisseurs~; APO10.05 surveiller la performance et la conformité des
fournisseurs.

\textbf{RACI (extrait, KRIBILOG)~:} APO10.05~: \emph{A} au responsable des
services informatiques, \emph{R} au gestionnaire du contrat, \emph{C} à la
direction des achats et au RSSI.

\textbf{Indicateurs~:} part des fournisseurs critiques évalués dans
l'année, nombre de manquements contractuels, délai de traitement des
écarts.
\end{boxobjet}
{\small\itshape Pratiques et objectifs liés traduits par nos soins à partir
de COBIT~2019~\cite{isaca2018cobitobjectifs}~; finalité reformulée.}

\subsection{Les clauses essentielles}

\begin{boxprocedure}
\textbf{Le contrat d'un fournisseur critique.} Pour qu'un auditeur puisse
établir que l'organisation maîtrise un service externalisé, le contrat
doit comporter au minimum~:
\begin{itemize}
    \item des engagements de service mesurables, dont au moins un
          indicateur de résultat, et la méthode de mesure~;
    \item des pénalités et un droit de résiliation pour manquements
          répétés~;
    \item l'obligation de signaler les incidents et les failles de sécurité
          dans un délai fixé~;
    \item un droit d'audit, par l'organisation ou par un tiers qu'elle
          désigne~;
    \item la localisation des données et les conditions de tout transfert
          (loi n°~2024/017, Chapitre~\ref{chap:droit})~;
    \item les conditions de la sous-traitance en cascade~;
    \item un plan de continuité du fournisseur, testé~;
    \item une clause de réversibilité~: restitution des données et
          assistance au transfert vers un autre prestataire.
\end{itemize}
\end{boxprocedure}

\begin{boxmetier}{Responsable des services IT}
\textbf{Le comité de service trimestriel avec NETPORT SERVICES}
\begin{enumerate}
    \item Revoir les indicateurs du trimestre, mesurés par KRIBILOG et non
          par le seul fournisseur.
    \item Examiner chaque incident P1 et P2, et les problèmes ouverts
          imputables au fournisseur.
    \item Arrêter les pénalités et les plans d'action.
    \item Revoir les risques du fournisseur (dépendance à une personne,
          santé financière, sous-traitants).
    \item Consigner les décisions au procès-verbal, transmis au comité
          stratégique des SI.
\end{enumerate}
\end{boxmetier}

\begin{boxmetier}{Auditeur SI}
\textbf{Auditer un fournisseur critique}
\begin{enumerate}
    \item Vérifier que le contrat contient les clauses essentielles
          ci-dessus.
    \item Recalculer, sur trois mois, la disponibilité et les pénalités à
          partir des données de l'organisation, et les comparer au rapport
          du fournisseur.
    \item Vérifier que les pénalités dues ont été effectivement appliquées.
    \item Vérifier le dernier test du plan de continuité du fournisseur.
\end{enumerate}
\end{boxmetier}

% ============================================================================
\separateur
% ============================================================================

\begin{boxresume}
\textbf{Ce qu'il faut retenir de ce chapitre~:}
\begin{itemize}
    \item \textbf{SLA, OLA, contrat de sous-traitance}~: on ne promet pas au
          client plus que ce que garantissent les accords qui soutiennent le
          service.
    \item \textbf{Un SLA utile} mesure au moins un résultat~; sinon
          apparaît l'effet pastèque.
    \item \textbf{La disponibilité se calcule}~: 99,5~\% sur 720~heures
          autorisent 3,6~heures d'interruption. En série, les
          disponibilités se multiplient~; en parallèle, les secours se
          composent.
    \item \textbf{Les pénalités} sont un signal, pas une assurance~: elles
          couvrent rarement le préjudice réel.
    \item \textbf{La continuité} part de l'analyse d'impact (DMIA, RTO,
          RPO) et ne vaut que si le plan est testé.
    \item \textbf{L'externalisation ne transfère pas la responsabilité}~:
          la loi n°~2010/012 vise les services fournis directement ou
          indirectement. Le contrat doit donner les moyens du contrôle.
\end{itemize}
\end{boxresume}

\bigskip

\begin{boxexo}[title={\ding{45}\quad Exercice~15.1. Calculs de disponibilité \niveaubase}]
\begin{enumerate}[(a)]
    \item Combien d'heures d'interruption par mois de 30~jours autorise une
          disponibilité de 99,8~\%~?
    \item Un service a connu 5~heures d'interruption en avril (30~jours).
          Quelle est sa disponibilité~?
    \item Un service dépend de trois composants en série disponibles à
          99,9~\%, 99,7~\% et 99,5~\%. Quelle est sa disponibilité~?
    \item On double le composant le plus faible par un secours identique.
          Quelle est la nouvelle disponibilité~?
\end{enumerate}
\end{boxexo}

\begin{boxexo}[title={\ding{45}\quad Exercice~15.2. Les pénalités de juillet \niveaumoyen}]
Avec la clause du scénario de la section~\ref{sec:c15-penalites}, le mois
de juillet (31~jours) a connu quatre interruptions de 3~heures,
2~heures, 6~heures et 1~heure.
\begin{enumerate}[(a)]
    \item Calculez la disponibilité du mois.
    \item Calculez la pénalité en FCFA.
    \item NETPORT SERVICES conteste l'interruption de 6~heures, qu'elle
          attribue à l'opérateur mobile. Que dit le contrat, et que
          devrait-il dire~?
\end{enumerate}
\end{boxexo}

\begin{boxexo}[title={\ding{45}\quad Exercice~15.3. La continuité de l'ANPP \niveauexpert}]
L'ANPP verse chaque mois les salaires de plus de 300~000 agents publics.
\begin{enumerate}[(a)]
    \item Conduisez une analyse d'impact simplifiée de ses trois services
          principaux (paie, paiement des fournisseurs, portail
          d'information), et fixez pour chacun DMIA, RTO et RPO.
    \item Proposez un dispositif de continuité proportionné, et le programme
          de tests de l'année.
    \item Identifiez les clauses à imposer à l'hébergeur de l'agence.
    \item Indiquez à quelle instance de gouvernance les résultats des tests
          doivent être présentés, et pourquoi.
\end{enumerate}
\end{boxexo}

\bigskip

\begin{boxinfo}
\textbf{Transition.} Les Chapitres~\ref{chap:itil-svs} à~\ref{chap:itil-niveaux}
ont présenté les pratiques d'ITIL~4 une par une. Le
Chapitre~\ref{chap:processus} apprend à les assembler~: modéliser un
processus, en attribuer les responsabilités par une matrice RACI, le
mesurer, et établir la correspondance entre COBIT, ITIL et
l'ISO/IEC~20000-1.
\end{boxinfo}
